\section{The Spin-Singlet Cooper Pair and Its Two Charge Channels}
\label{sec:ch21-singlet-pair}

\subsection*{Why one singlet operator has two different targets}

A spin-singlet Cooper pair is often represented by a single compact field expression in the continuum literature. However, on a finite fermionic Hilbert space, we must first ask a more elementary question: \emph{which occupation-number sectors does this operator connect?} When assembling a pair from a right mover and a left mover, the two antisymmetrized spin assignments remove particles from different species. Although they carry the same \emph{total} charge and spin, they do not possess the same four-component charge vector. Consequently, when their finite energy budgets are treated as typed linear maps, their operators have distinctly different target spaces. Furthermore, their relative fermionic signs must be explicitly computed and cannot simply be discarded under the heuristic phrase ``the Klein factors cancel.''

We continue with the ordered four-species CAR carrier of
Section~\ref{sec:ch21-spin-charge}.  Set
\begin{equation}
 L=2h,\quad h\in\mathbb N_{>0},\qquad
 \mathcal B_L=\{1-h,2-h,\ldots,h\},\qquad
 \Lambda_L=\mathbb Z/L\mathbb Z,\qquad
 \zeta=e^{2\pi i/L}.
 \label{eq:ch21-pair-setting}
\end{equation}
Every one of the four species
\(\mathcal C=\{L,R\}\times\{\uparrow,\downarrow\}\) has these \(L\)
momentum modes; the lexicographic occupation order is
\begin{equation}
 L\uparrow<L\downarrow<R\uparrow<R\downarrow,
 \quad\text{then increasing momentum within each species}.
 \label{eq:ch21-pair-order}
\end{equation}
We retain the positive-mode cutoff notation \(1\leq m\leq M\) for later
current or oscillator expansions, but the exact CAR identities in this
section require no such cutoff.  All site labels \(x\) are periodic lattice
positions, represented by integers modulo \(L\).  With periodic Fourier
orientation fixed as in Section~\ref{sec:ch21-spin-charge}, the finite site
annihilator is
\begin{equation}
 c_{\alpha,x}=\frac1{\sqrt L}\sum_{k\in\mathcal B_L}
        \zeta^{kx}c_{\alpha,k},\qquad
 c_{\alpha,x}^{\dagger}=\frac1{\sqrt L}\sum_{k\in\mathcal B_L}
        \zeta^{-kx}c_{\alpha,k}^{\dagger}.
 \label{eq:ch21-pair-site-field}
\end{equation}
These are genuine operators on the finite occupation Fock space
\(\mathcal H_L\), not formal vertex exponentials.  Orthogonality of the
finite characters and the momentum CAR give
\(\{c_{\alpha,x},c_{\beta,y}^{\dagger}\}
=\delta_{\alpha\beta}\delta_{xy}I\) and
\(\{c_{\alpha,x},c_{\beta,y}\}=0\).
Indeed, the first anticommutator is
\(L^{-1}\delta_{\alpha\beta}\sum_{k\in\mathcal B_L}
\zeta^{k(x-y)}I\), which is the stated periodic Kronecker delta.
There is no non-wrapping current approximation in this step.

\begin{definition}[Annihilating spin-singlet pair]
\label{def:ch21-singlet}
The ordered, spin-antisymmetric two-fermion operator is
\begin{equation}
 \begin{split}
 O_{SSC}(x)&=A_1(x)-A_2(x),\\
 A_1(x)&=c_{R\uparrow,x}\,c_{L\downarrow,x},\\
 A_2(x)&=c_{R\downarrow,x}\,c_{L\uparrow,x}.
 \end{split}
 \label{eq:ch21-singlet-def}
\end{equation}
The rightmost field acts first.  In particular, \(A_1\) annihilates an
\(L\downarrow\) particle before an \(R\uparrow\) particle; \(A_2\)
annihilates an \(L\uparrow\) particle before an \(R\downarrow\) particle.
\end{definition}

The minus sign in \eqref{eq:ch21-singlet-def} is the antisymmetric spin
contraction \(\sum_{\sigma,\tau}\epsilon_{\sigma\tau}
 c_{R\sigma,x}c_{L\tau,x}\), with
\(\epsilon_{\uparrow\downarrow}=1\) and
\(\epsilon_{\downarrow\uparrow}=-1\).  For \(U\in SU(2)\), the elementary
identity \(U^{T}\epsilon U=(\det U)\epsilon=\epsilon\) shows that this
contraction is a spin scalar when both spin labels transform by the same
matrix.  This representation-theoretic assertion does \emph{not} say that
the Hamiltonian possesses spin-rotation symmetry.  It also does not
eliminate the antisymmetric CAR ordering.  For example, taking the adjoint
and reordering the fermions gives
\begin{equation}
 O_{SSC}(x)^\dagger
  =-c_{R\uparrow,x}^\dagger c_{L\downarrow,x}^\dagger
    -c_{L\uparrow,x}^\dagger c_{R\downarrow,x}^\dagger.
 \label{eq:ch21-singlet-adjoint}
\end{equation}
Thus the annihilating convention here differs by an overall sign upon
adjunction from a frequently written spin-singlet \emph{creation} convention.

\subsection*{Four-component charge shifts and the common total quantum numbers}

Let \(\widehat N_\alpha=\sum_{k\in\mathcal B_L}
 c_{\alpha,k}^\dagger c_{\alpha,k}-hI\) be the relative species charge;
write \(\mathbf N\in\mathbb Z^{\mathcal C}\) for its simultaneous
eigenvalues, and \(e_\alpha\) for the unit charge vector in species
\(\alpha\).  Only admissible components \(-h\leq N_\alpha\leq h\)
occur on the finite carrier.  The \(c_{\alpha,k}\) preserve neither
\(\mathbf N\) nor the energy budget, but their \emph{charge shifts} are
exact on the entire ambient finite Fock space.

\begin{lemma}[Exact species shifts of the two singlet channels]
\label{lem:ch21-singlet-shifts}
For every \(\gamma\in\mathcal C\) and every \(x\in\Lambda_L\),
\begin{align}
 [\widehat N_\gamma,A_1(x)]
   &=-(\delta_{\gamma,R\uparrow}+\delta_{\gamma,L\downarrow})A_1(x),
   \label{eq:ch21-pair-shift-one}\\
 [\widehat N_\gamma,A_2(x)]
   &=-(\delta_{\gamma,R\downarrow}+\delta_{\gamma,L\uparrow})A_2(x).
   \label{eq:ch21-pair-shift-two}
\end{align}
Consequently, if \(\mathcal H_{\mathbf N}\subseteq\mathcal H_L\)
is the full joint eigenspace of species charges \(\mathbf N\), then
\begin{equation}
 A_j(x)\mathcal H_{\mathbf N}\subseteq\mathcal H_{\mathbf T_j},
 \qquad\begin{cases}
 \mathbf T_1=\mathbf N-e_{L\downarrow}-e_{R\uparrow},\\
 \mathbf T_2=\mathbf N-e_{L\uparrow}-e_{R\downarrow}.
 \end{cases}
 \label{eq:ch21-pair-targets}
\end{equation}
These two targets are distinct integer vectors.  Their charge eigenspaces
are orthogonal, including when one of the eigenspaces is zero.
\end{lemma}

\begin{proof}
The CAR imply
\([c_{\gamma,k}^\dagger c_{\gamma,k},c_{\beta,q}]
=-\delta_{\gamma\beta}\delta_{kq}c_{\beta,q}\).
Summing over \(k\), subtracting the scalar sea occupation \(hI\), and
then using \eqref{eq:ch21-pair-site-field} gives
\([\widehat N_\gamma,c_{\beta,x}]
=-\delta_{\gamma\beta}c_{\beta,x}\).
For any operators \(D,X,Y\), the identity
\([D,XY]=[D,X]Y+X[D,Y]\) holds without a commutativity
assumption.  Apply it to the \emph{ordered} words \(A_1,A_2\) to obtain
\eqref{eq:ch21-pair-shift-one}--\eqref{eq:ch21-pair-shift-two}.
If \(\widehat N_\gamma\psi=N_\gamma\psi\), then
\(\widehat N_\gamma A_j\psi=(N_\gamma+\Delta_{j,\gamma})A_j\psi\),
which gives \eqref{eq:ch21-pair-targets}.  Its two shift vectors are
different because, for instance, their \(L\uparrow\) components are
\(0\) and \(-1\).  Finally, distinct eigenspaces of a self-adjoint
operator are orthogonal: if \(\widehat N_\gamma u=au\) and
\(\widehat N_\gamma v=bv\), then
\(a\langle u,v\rangle=\langle\widehat N_\gamma u,v\rangle
=\langle u,\widehat N_\gamma v\rangle=b\langle u,v\rangle\).
Choose a component for which \(a\ne b\).
\end{proof}

It is instructive to express these shifts using the raw charge and spin
labels of Section~\ref{sec:ch21-spin-charge}:
\(Q_{\eta,c}=N_{\eta\uparrow}+N_{\eta\downarrow}\) and
\(Q_{\eta,s}=N_{\eta\uparrow}-N_{\eta\downarrow}\).
\begin{center}
\begin{tabular}{@{}lrrrrrr@{}}
\toprule
 & $\Delta Q_{R,c}$ & $\Delta Q_{R,s}$
 & $\Delta Q_{L,c}$ & $\Delta Q_{L,s}$
 & $\Delta Q_{\mathrm{tot},c}$ & $\Delta Q_{\mathrm{tot},s}$\\
\midrule
$A_1$ & $-1$ & $-1$ & $-1$ & $+1$ & $-2$ & $0$\\
$A_2$ & $-1$ & $+1$ & $-1$ & $-1$ & $-2$ & $0$\\
\bottomrule
\end{tabular}
\end{center}
The common total \(\Delta Q_{\mathrm{tot},c}=-2\) and
\(\Delta Q_{\mathrm{tot},s}=0\) do not identify the two targets.  In
particular, the chirality-resolved spin shifts differ, and each branch
continues to obey the integer parity constraint
\(Q_{\eta,c}\equiv Q_{\eta,s}\pmod 2\).  Forgetting chirality-resolved
charges before projecting would erase information needed to determine the
correct codomain.

\subsection*{The signed Klein factors are ordered charge-shift maps}

While a Klein factor is sometimes introduced as an abstract anticommuting symbol attached to each fermion field, on our finite occupation representation the source of its sign is entirely concrete: removing a fermion requires commuting it past all occupied modes that precede it in the fixed species--momentum order. Thus, the signed factor is a rigorously defined \emph{charge-dependent map}, rather than an arbitrary numerical phase assigned loosely to each species.

For each admissible integer vector \(\mathbf N\), write
\(\Omega_{\mathbf N}\) for the normalized charge-sector ground ket, with
species \(\alpha\) filled at
\(\{1-h,\ldots,N_\alpha\}\).  Its number of particles of species
\(\alpha\) is \(h+N_\alpha\).  Assume below that a named annihilation is
possible, namely \(N_\alpha\geq1-h\), and that all indicated lowered
charge vectors remain admissible.  In Chapter~\ref{ch:car_fock}'s
occupation action, the annihilator of the \emph{top occupied momentum}
\(k=N_\alpha\) has preceding-occupation exponent
\begin{equation}
 \begin{split}
 E_h(\mathbf N,\alpha)
   &=\sum_{\beta<\alpha}(h+N_\beta)
       +(h+N_\alpha-1),\\
 \kappa_h(\mathbf N,\alpha)
   &=(-1)^{E_h(\mathbf N,\alpha)}\in\{+1,-1\}.
 \end{split}
 \label{eq:ch21-pair-klein-phase}
\end{equation}
The first term counts occupied modes in preceding species; the second
counts those below the occupied top mode within species \(\alpha\).
The signed finite Klein kernel constructed in Core Chapter~13
reproduces precisely this phase on sector ground states, under the
charge and budget margins stated for that construction.

To keep the zero-mode notation transparent, denote by \(F_\alpha\) the
\emph{typed sector shift} whose ground-state action is
\begin{equation}
 F_\alpha\Omega_{\mathbf N}
  =\kappa_h(\mathbf N,\alpha)\,
    \Omega_{\mathbf N-e_\alpha}.
 \label{eq:ch21-pair-klein-action}
\end{equation}
This notation describes the Core Klein action on admissible ground
states; extending it to particular non-ground budgets requires the
corresponding finite kernel theorem.  It is not an assertion that one
bounded lowering operator acts isometrically on every charge sector,
including the empty boundary sector.  Define also the diagonal site phase
\(Z_{\alpha,x}\Omega_{\mathbf N}=\zeta^{xN_\alpha}
\Omega_{\mathbf N}\).  Since lowering species \(\beta\ne\alpha\)
does not change \(N_\alpha\), the ground-state actions of
\(Z_{\alpha,x}\) and \(F_\beta\) commute for distinct species.

\begin{proposition}[Two different signed zero-mode products]
\label{prop:ch21-two-klein-products}
With \(\mathbf T_1,\mathbf T_2\) as in
\eqref{eq:ch21-pair-targets}, the two ordered zero-mode combinations are
\begin{align}
 \mathcal K_1(x)&=
 F_{R\uparrow}F_{L\downarrow}
 Z_{R\uparrow,x}Z_{L\downarrow,x}:
  \mathbf N\longrightarrow\mathbf T_1,
 \label{eq:ch21-klein-one}\\
 \mathcal K_2(x)&=
 F_{R\downarrow}F_{L\uparrow}
 Z_{R\downarrow,x}Z_{L\uparrow,x}:
  \mathbf N\longrightarrow\mathbf T_2.
 \label{eq:ch21-klein-two}
\end{align}
On an admissible sector ground ket, the actions are
\begin{equation}
 \mathcal K_j(x)\Omega_{\mathbf N}
  =\zeta^{x(N_{a_j}+N_{b_j})}
    \kappa_h(\mathbf N-e_{b_j},a_j)
    \kappa_h(\mathbf N,b_j)\,\Omega_{\mathbf T_j},
 \label{eq:ch21-klein-channel-ground}
\end{equation}
where \((a_1,b_1)=(R\uparrow,L\downarrow)\) and
\((a_2,b_2)=(R\downarrow,L\uparrow)\).
The rightmost \(F_{b_j}\) acts first; in particular the second
Klein phase must be evaluated on \(\mathbf N-e_{b_j}\), not on
\(\mathbf N\).
\end{proposition}

\begin{proof}
Act with the rightmost diagonal \(Z\) operators on
\(\Omega_{\mathbf N}\); their product is
\(\zeta^{x(N_{a_j}+N_{b_j})}\).  Next \(F_{b_j}\) lowers
\(\mathbf N\) to \(\mathbf N-e_{b_j}\) with phase
\(\kappa_h(\mathbf N,b_j)\).  Finally \(F_{a_j}\) acts on that
\emph{intermediate} charge vector, giving
\(\kappa_h(\mathbf N-e_{b_j},a_j)\).
This proves \eqref{eq:ch21-klein-channel-ground}.
Notice that in our fixed ordering both \(b_1<a_1\) and
\(b_2<a_2\).  Hence
\(E_h(\mathbf N-e_b,a)=E_h(\mathbf N,a)-1\) whenever
\(b<a\) and \(a\ne b\).  Reversing the two shifts changes their
relative sign, which is the finite occupation-string origin of the
fermionic cross-species anticommutation.  No freedom to suppress that
sign has been introduced.
\end{proof}

The \(Z\) factors specify zero-mode phases, not the full oscillator
content of the site fields.  The formulas above are \emph{not} a
factorization theorem expressing \(c_{a,x}c_{b,x}\) on every budget
as a product of two compressed bosonic exponentials.  Such a statement
requires the single-field identifications, a common domain of composition,
and control of intermediate projection remainders.

\begin{lemma}[Matching the ordered CAR action on sector grounds]
\label{lem:ch21-pair-ground-action}
For distinct species \(a,b\), assume that the charge-sector ground
\(\Omega_{\mathbf N}\) and both lowerings are admissible.  Let
\(k_b=N_b\) and \(k_a=N_a\) denote the occupied top momenta in
\(\mathcal B_L\).  Then
\begin{equation}
 c_{a,k_a}c_{b,k_b}\Omega_{\mathbf N}
  =\kappa_h(\mathbf N-e_b,a)\kappa_h(\mathbf N,b)
      \Omega_{\mathbf N-e_b-e_a},
 \label{eq:ch21-pair-ordered-top}
\end{equation}
and, for the site fields,
\begin{equation}
 \big\langle\Omega_{\mathbf N-e_b-e_a},
  c_{a,x}c_{b,x}\Omega_{\mathbf N}\big\rangle
   =\frac{\zeta^{x(N_a+N_b)}}{L}
     \kappa_h(\mathbf N-e_b,a)\kappa_h(\mathbf N,b).
 \label{eq:ch21-pair-channel-coefficient}
\end{equation}
The right-hand side is \(L^{-1}\) times the corresponding
zero-mode matrix coefficient in \eqref{eq:ch21-klein-channel-ground}.
\end{lemma}

\begin{proof}
For a ground configuration, deleting its top occupied momentum
\(k_b=N_b\) leaves the ground configuration of charge
\(\mathbf N-e_b\); the occupation-basis annihilation rule gives the
phase \(\kappa_h(\mathbf N,b)\).  Deleting \(k_a=N_a\) from this
intermediate ground gives the second phase
\(\kappa_h(\mathbf N-e_b,a)\), proving
\eqref{eq:ch21-pair-ordered-top}.
To compute the matrix element for site annihilators, expand both
finite Fourier sums in \eqref{eq:ch21-pair-site-field}.
The target ket is the \emph{ground} configuration at
\(\mathbf N-e_a-e_b\).  Removing any momentum other than the
occupied top mode of species \(b\) leaves a different occupation
configuration, and removal of a non-top momentum of species \(a\)
likewise cannot produce the target ground.  Orthogonality of
occupation kets therefore selects the unique summand
\((k_a,k_b)=(N_a,N_b)\).  Its two Fourier amplitudes multiply to
\(L^{-1}\zeta^{x(N_a+N_b)}\), while its CAR sign is the product
already obtained.  No current-CCR or energy-margin hypothesis enters
this matrix coefficient; it is an exact finite-CAR statement.
\end{proof}

\subsection*{A finite sign calculation: both Klein phases matter}

\begin{example}[Four species, \(L=4\), and two opposite singlet coefficients]
\label{ex:ch21-pair-sign}
Take \(h=2\), \(L=4\), \(x=1\), \(\zeta=i\), and the zero-charge sea
\(\Omega_{\mathbf0}\).  Each of the four species occupies momenta
\(\{-1,0\}\), ordered as in \eqref{eq:ch21-pair-order}.
Both singlet channels remove the top momentum \(k=0\) from their
respective species.

For channel 1, annihilating the rightmost \(L\downarrow\) field
encounters \(2+1=3\) preceding occupied modes, hence contributes
\(-1\).  In the intermediate state, annihilating \(R\uparrow\) at
\(k=0\) encounters \(2+1+1=4\) preceding modes, hence contributes
\(+1\).  Thus
\begin{equation}
 \kappa_2(\mathbf0,L\downarrow)=-1,\quad
 \kappa_2(-e_{L\downarrow},R\uparrow)=+1,\quad
 \langle\Omega_{\mathbf T_1},A_1(1)\Omega_{\mathbf0}\rangle=-\tfrac14.
 \label{eq:ch21-pair-example-one}
\end{equation}
For channel 2, removing \(L\uparrow\) at \(k=0\) encounters one
preceding occupied mode and contributes \(-1\).  Removing
\(R\downarrow\) at \(k=0\) from the intermediate configuration
encounters \(1+2+2+1=6\) preceding modes and contributes \(+1\).
Therefore
\begin{equation}
 \kappa_2(\mathbf0,L\uparrow)=-1,\quad
 \kappa_2(-e_{L\uparrow},R\downarrow)=+1,\quad
 \langle\Omega_{\mathbf T_2},A_2(1)\Omega_{\mathbf0}\rangle=-\tfrac14.
 \label{eq:ch21-pair-example-two}
\end{equation}
The Fourier phase is \(\zeta^{(0+0)x}=1\) in both cases, and each
pair contributes the normalization \(1/L=1/4\).  In the \emph{singlet
operator} the second channel is subtracted.  Accordingly,
\begin{equation}
 \boxed{\quad
 \langle\Omega_{\mathbf T_1},O_{SSC}(1)\Omega_{\mathbf0}\rangle
   =-\tfrac14,\qquad
 \langle\Omega_{\mathbf T_2},O_{SSC}(1)\Omega_{\mathbf0}\rangle
   =+\tfrac14.\quad}
 \label{eq:ch21-pair-example-final}
\end{equation}
There is no contradiction in the opposite signs: the target kets
belong to different, mutually orthogonal charge sectors.  Their
amplitudes must not be summed as if they were two contributions to
the same basis coefficient.
\end{example}

The sign pattern has a useful exact generalization.  With
\(\mathbf N=\mathbf0\) and arbitrary \(h\geq1\), the channel-1
preceding exponents are \(2h-1\) for \(L\downarrow\) and
\(3h-2\) for \(R\uparrow\) after the first removal.  Channel 2 has
exponents \(h-1\) for \(L\uparrow\) and \(4h-2\) for
\(R\downarrow\) after the first removal.  Therefore
\begin{equation}
 \kappa_h(-e_{L\downarrow},R\uparrow)
 \kappa_h(\mathbf0,L\downarrow)
 =\kappa_h(-e_{L\uparrow},R\downarrow)
 \kappa_h(\mathbf0,L\uparrow)=(-1)^{h+1}.
 \label{eq:ch21-pair-general-sign}
\end{equation}
At \(\mathbf N=\mathbf0\) the ground-to-ground coefficients are
\(\pm(-1)^{h+1}/L\) in the two mutually orthogonal targets,
with the overall negative sign applied to channel 2.

\subsection*{Typed finite budgets, the direct-sum target, and exact residuals}

The finite-field product is an ambient endomorphism on \(\mathcal H_L\), meaning \(A_1-A_2\) is well defined there. Trouble arises only if each annihilator is first compressed to a map between \emph{different} charge-sector energy budgets and these maps are then subtracted as if their codomains agreed. The remedy is to rigorously specify each intermediate space, calculate its omitted component, and embed both final targets into the same encompassing Hilbert space.

Let \(\mathcal E(S)\) be the nonnegative integer grading relative to
the ground configuration of the charge sector of \(S\), as in
Chapter~\ref{ch:vacuum_budget}.  For admissible \(\mathbf N\) and
\(J\in\mathbb N\), define the coordinate budget
\begin{equation}
 V_{\mathbf N,J}
  =\operatorname{span}\{\,|S\rangle:\mathbf N(S)=\mathbf N,
                          \ \mathcal E(S)\leq J\,\}
  \subseteq\mathcal H_L.
 \label{eq:ch21-pair-budget}
\end{equation}
Let \(i_{\mathbf N,J}:V_{\mathbf N,J}\hookrightarrow\mathcal H_L\)
be inclusion, \(p_{\mathbf N,J}=i_{\mathbf N,J}^\dagger\) its
orthogonal component projection, and
\(P_{\mathbf N,J}=i_{\mathbf N,J}p_{\mathbf N,J}\) the ambient
orthogonal projector.  A cutoff \(J\) is an excitation-energy budget,
not a bound on the scalar product \(\langle\psi,H\psi\rangle\).

For \(j=1,2\), retain the ordered pair \((a_j,b_j)\) above and choose
\emph{independent} source, intermediate, and target cutoffs
\(J_s,J_{m,j},J_{t,j}\in\mathbb N\).  Define
\begin{equation}
 \begin{aligned}
 V_s&=V_{\mathbf N,J_s},\\
 \mathbf M_j&=\mathbf N-e_{b_j},
    &V_{m,j}&=V_{\mathbf M_j,J_{m,j}},\\
 \mathbf T_j&=\mathbf N-e_{b_j}-e_{a_j},
    &V_{t,j}&=V_{\mathbf T_j,J_{t,j}}.
 \end{aligned}
 \label{eq:ch21-pair-three-spaces}
\end{equation}
If any shifted charge vector is not admissible, its corresponding
budget is interpreted as the zero subspace of \(\mathcal H_L\).
Write \(i_s,p_{m,j},i_{m,j},p_{t,j}\) for the associated inclusion
and coordinate projections.  Finally form the \emph{orthogonal sum}
\begin{equation}
 W=V_{t,1}\oplus V_{t,2},\qquad
 j_j:V_{t,j}\longrightarrow W,\qquad
 p_W:\mathcal H_L\longrightarrow W.
 \label{eq:ch21-pair-common-target}
\end{equation}
Here \(W\) is identified isometrically with the subspace
\(V_{t,1}+V_{t,2}\subseteq\mathcal H_L\); the sum is orthogonal by
Lemma~\ref{lem:ch21-singlet-shifts}.  Thus each \(j_j\) is the
canonical inclusion of one target into the common codomain.

\begin{theorem}[Exact two-channel projected singlet identity]
\label{thm:ch21-pair-projected}
For every finite \(L=2h\), site \(x\), admissible source charge,
and chosen nonnegative cutoffs, define the typed field maps
\begin{equation}
 \begin{aligned}
 B_{b,j}&=p_{m,j}c_{b_j,x}i_s:
                  V_s\longrightarrow V_{m,j},\\
 B_{a,j}&=p_{t,j}c_{a_j,x}i_{m,j}:
                  V_{m,j}\longrightarrow V_{t,j},\\
 \mathcal R_j&=p_{t,j}c_{a_j,x}
       (I-P_{m,j})c_{b_j,x}i_s:
                  V_s\longrightarrow V_{t,j},\\
 P_{m,j}&=i_{m,j}p_{m,j}:
                  \mathcal H_L\longrightarrow\mathcal H_L.
 \end{aligned}
 \label{eq:ch21-pair-typed-maps}
\end{equation}
Then, as an equality of linear maps \(V_s\longrightarrow W\),
\begin{equation}
 \boxed{\displaystyle
 p_W O_{SSC}(x)i_s
  =j_1\!\left(B_{a,1}\circ B_{b,1}+\mathcal R_1\right)
   -j_2\!\left(B_{a,2}\circ B_{b,2}+\mathcal R_2\right).}
 \label{eq:ch21-pair-typed-identity}
\end{equation}
No intermediate or final cutoff inequality is required.  In
particular, \(\mathcal R_1\) and \(\mathcal R_2\) are separately
specified, and neither is silently set to zero.
\end{theorem}

\begin{proof}
Fix a channel \(j\).  The insertion of the ambient orthogonal
identity \(I=P_{m,j}+(I-P_{m,j})\) \emph{between the two site
annihilators} is exact:
\begin{equation}
 \begin{aligned}
 p_{t,j}c_{a_j,x}c_{b_j,x}i_s
 &=p_{t,j}c_{a_j,x}
       P_{m,j}c_{b_j,x}i_s
   +p_{t,j}c_{a_j,x}(I-P_{m,j})c_{b_j,x}i_s\\
 &=\bigl(p_{t,j}c_{a_j,x}i_{m,j}\bigr)
       \bigl(p_{m,j}c_{b_j,x}i_s\bigr)
      +\mathcal R_j\\
 &=B_{a,j}\circ B_{b,j}+\mathcal R_j.
 \end{aligned}
 \label{eq:ch21-pair-insert-projection}
\end{equation}
This equality holds inside \(\operatorname{Hom}(V_s,V_{t,j})\).
By Lemma~\ref{lem:ch21-singlet-shifts}, the image of
\(c_{a_j,x}c_{b_j,x}i_s\) lies entirely in the
\emph{full} charge sector \(\mathcal H_{\mathbf T_j}\).
Since \(V_{t,1}\) and \(V_{t,2}\) occupy different charge sectors,
projecting the channel-\(j\) image to \(W\) yields precisely
\(j_jp_{t,j}c_{a_j,x}c_{b_j,x}i_s\); the projection to the other
target is zero.  Inserting the two channel identities
\eqref{eq:ch21-pair-insert-projection} into
\(p_W(A_1-A_2)i_s\) proves
\eqref{eq:ch21-pair-typed-identity}.  The two terms are subtracted
\emph{after} composition with \(j_1\) and \(j_2\), so their
codomain is genuinely the same \(W\).  This completes the proof.
\end{proof}

The theorem is a finite-Fock identity, not an approximate bosonization
formula.  Both residuals represent concrete vectors that passed through
the \emph{complement} of the selected intermediate energy budget before
being annihilated by the second field.  In particular,
\begin{equation}
 \mathcal R_j=0
 \quad\Longleftrightarrow\quad
 p_{t,j}c_{a_j,x}(I-P_{m,j})c_{b_j,x}i_s=0.
 \label{eq:ch21-pair-residual-criterion}
\end{equation}
The stronger range condition
\((I-P_{m,j})c_{b_j,x}i_s=0\) suffices, but is not necessary:
the second field and final projection can remove all discarded
components even when the intermediate field leaves the budget.  A
single-field bosonization candidate from Core Chapter~14 may replace
\(B_{a,j}\) or \(B_{b,j}\) only after its \emph{individual}
finite-budget identity and margins have been established.  A verified product of independently
compressed candidates is not automatically the compressed physical
product: precisely the \(\mathcal R_j\) term measures the difference.

\begin{example}[A nonzero intermediate-projection residual]
\label{ex:ch21-pair-residual}
Return to \(h=2\), \(L=4\), \(x=1\), source charge
\(\mathbf N=\mathbf0\) and source budget \(J_s=0\).  In channel 1,
take \(J_{m,1}=0\) but \(J_{t,1}=1\).  Starting from the sea
\(\Omega_{\mathbf0}\), remove momentum \(k=-1\), rather than the
top momentum \(k=0\), in species \(L\downarrow\).  The resulting
species occupation is \(\{0\}\), whereas the ground occupation of
relative charge \(-1\) is \(\{-1\}\).  It therefore has one unit
of excitation energy and is annihilated by the intermediate projection
\(P_{m,1}\).

Next remove the top momentum \(k=0\) in species \(R\uparrow\).
The final occupation ket \(\Xi\) has charge
\(\mathbf T_1=-e_{L\downarrow}-e_{R\uparrow}\), one excitation
unit in \(L\downarrow\), and zero excitation in every other
species, so \(\Xi\in V_{t,1}\).  No other momentum pair
produces this exact configuration.  The preceding counts are \(2\)
for the first \(L\downarrow\) removal and \(4\) for the
\(R\uparrow\) removal; both CAR signs are positive.  The two
Fourier factors contribute
\(L^{-1}\zeta^{-1}=(-i)/4\).  Consequently,
\begin{equation}
 \left\langle\Xi,\mathcal R_1\Omega_{\mathbf0}\right\rangle
     =-\frac{i}{4}\ne0,\qquad
 \left\langle\Xi,(B_{a,1}\circ B_{b,1})
          \Omega_{\mathbf0}\right\rangle=0.
 \label{eq:ch21-pair-residual-witness}
\end{equation}
Thus even a source \emph{ground state} does not make a truncated
intermediate projection exact for an arbitrary final cutoff.  Taking
\(J_{m,1}=1\) retains the displayed intermediate ket, but that alone
is not a proof that the entire residual vanishes for all source
vectors and all final cutoffs.
\end{example}

\begin{remark}[An enlarged final codomain]
\label{rem:ch21-pair-enlarged-target}
Equation~\eqref{eq:ch21-pair-typed-identity} uses \(W\) \emph{equal to}
the indicated two-budget orthogonal sum.  If \(W'\) is any larger
closed subspace containing \(W\), with inclusion
\(\iota:W\hookrightarrow W'\), its orthogonal projection satisfies
\begin{equation}
 p_{W'}O_{SSC}(x)i_s
  =\iota\,p_WO_{SSC}(x)i_s
     +p_{W'}(I-P_W)O_{SSC}(x)i_s.
 \label{eq:ch21-pair-larger-target}
\end{equation}
The last term includes whatever part of either channel lies in the
larger target outside the two selected budgets; it must be retained.
It vanishes if \(W'=W\), or under an independently proved range
condition.  Thus merely saying that \(W'\) ``contains both target
budgets'' is insufficient to infer the same two-term formula.
\end{remark}

\subsection*{What is established, and what remains conditional}

In summary, we have derived three logically separate statements.  First,
\eqref{eq:ch21-pair-shift-one}--\eqref{eq:ch21-pair-targets} are exact
species-charge selection rules in the full finite CAR algebra.  They
show why the singlet has two, rather than one, typed charge channels.
Second, \eqref{eq:ch21-pair-klein-phase}--\eqref{eq:ch21-pair-general-sign}
compute their ordered fermionic phases on charge-sector ground kets;
these phases are the ground-state content of the finite signed Klein
construction.  Third, \eqref{eq:ch21-pair-typed-identity} is an exact
linear-map identity on arbitrary finite coordinate budgets, and
Example~\ref{ex:ch21-pair-residual} shows why its two residuals cannot
in general be omitted.

In a charge/spin collective-field description, one may substitute the
algebraic inverse currents
\(\rho_{\eta\uparrow}=(R^c_\eta+R^s_\eta)/2\) and
\(\rho_{\eta\downarrow}=(R^c_\eta-R^s_\eta)/2\) into candidate
oscillator phases.  That substitution is \emph{not} by itself a proof
that compressed exponentials factor into a single common charge factor
and a common spin factor.  The two singlet channels have different
Klein products, different charge-sector targets, and potentially
different intermediate-budget errors.  Any such rewriting must specify
the uncompressed algebra in which exponentials are defined, prove the
required commutation and ordering identities, and establish their
compatibility with projection.  The formal series of
Section~\ref{sec:ch20-formal-vertices} likewise cannot be evaluated
at a numerical parameter or identified with finite-CAR expectations
without further arguments.

No thermodynamic-limit correlation exponent, spectral gap, or
interacting-state identification follows from the operator identities
of this section.  The exact CAR results used here rest on the
previously constructed finite carrier and its supported Klein and
coordinate-projection results; the spinful Chapter~21 implementations
and the present manuscript remain to be independently formalized,
compiled against the full book, and reviewed.  In the next section,
Umklapp scattering changes four species charges at once, making the
same sector-typing and projection questions essential rather than
merely notational.
